% Source pp. 1-9 (Lectures1-6.pdf).
\section*{Introduction}
\addcontentsline{toc}{section}{Introduction}

The subject of Field Theory is the dynamical systems (classical or quantum)
with continuously many degrees of freedom. That is, there is some number of
degrees of freedom associated with each point of space $\xx\in\R^3$. Examples
are:
\begin{enumerate}[label=\roman*)]
\item Electromagnetic Theory
\[
\text{Degrees of Freedom:}\qquad \mathbf{E}(\xx,t),\ \mathbf{H}(\xx,t)\,.
\]
\item Gravitation (general Relativity)
\[
\text{Degrees of Freedom} = \text{Spacetime metric}\quad g_{\mu\nu}(\xx,t)\,.
\]
\item More exotic example: Dirac (electron-positron) field
\[
\text{Degrees of Freedom} = \text{Four-component (complex) Lorentz spinor}\quad \psi_a(\xx,t)\,.
\]
\end{enumerate}

Main subject of this course is Quantum Theory of Fields. The states of
quantum fields with local interaction admit interpretation in terms of
propagating and interacting particles. Thus the theory of quantum fields is
the basis of particle theory.

\section{Elements of Classical Field Theory}

We start with a brief review of basic aspects of classical field theory.
Fundamental object in classical mechanics is the \textbf{action},
\[
S=\int L\bigl(q(t),\dot q(t)\bigr)\,dt\,;
\]
the equations of motion are derived from the extremal action principle
$\delta S=0$. In field theory the coordinates are replaced by the field
degrees of freedom,
\[
q(t)=\{q^i(t)\}\to \Phi(\xx,t)\,,\qquad (i\to\xx)
\]
and the action is written as a space-time integral
\[
S=\int \Lag\bigl(\Phi(\xx,t),\dot\Phi(\xx,t),\nabla\Phi(\xx,t)\bigr)\,d^3\xx\,dt\,.
\]
Here $\Phi=\Phi(\xx,t)$ stands for a collection of fields (functions of
$\xx,t$, representing local degrees of freedom), which may include scalar
fields, components of vectors, tensors, spinors, etc. It is conventionally
assumed that $\Phi$ has finitely many individual components. The function
$\Lag$ is called the \textbf{Lagrangian density}; it depends on $\Phi(\xx,t)$
and its derivatives over $\xx$ and $t$, taken at the same space-time point
$(\xx,t)$. Two remarks are in order.

1. The above expression is suitable for imposing conditions of the special
relativity. The integration is over the 4-dimensional Minkowskian space-time,
\[
d^3\xx\,dt=d^4x\,.
\]
Recall that in relativistic theory one deals with 4-dimensional space-time
with the coordinates
\[
x^\mu=(x^0,x^1,x^2,x^3)=(t,\xx)\,,
\]
equipped with the ``pseudo-metric''
\[
ds^2=-c^2\,dt^2+(dx^1)^2+(dx^2)^2+(dx^3)^2\,.
\]
The Lorentz transformations are those linear transformations of the
coordinates $x^\mu$ which preserve $ds^2$. In what follows I will use the
units in which\footnote{These conditions are simply the convention to measure
spatial distances and energy in the units
$[\text{unit of length}]=c\times[\text{unit of time}]$,
$[\text{unit of energy}]=\hbar/[\text{unit of time}]$.}
\[
c=1\,,\qquad \hbar=1\,.
\]
Then
\[
ds^2=\eta_{\mu\nu}\,dx^\mu dx^\nu=dx^\mu\,dx_\mu\,,\qquad\text{where}\qquad x^\mu=(t,\xx)\,,\quad x_\mu=(-t,\xx)\,,
\]
and $\eta_{\mu\nu}=\operatorname{diag}(-1,+1,+1,+1)$ is the Minkowski metric.
Throughout these notes the metric is ``mostly plus'', so a timelike vector
has $v^2=v^\mu v_\mu<0$ and an on-shell momentum has $p^2=-m^2$. In this
convention $px=-p^0x^0+\pp\cdot\xx$.
To ensure relativistic invariance of the classical field theory it suffices
to choose $\Lag$ to be a Lorentz scalar.

2. In writing the above action I already accepted severe restrictions,
namely,
\[
S=\int dt\,L\,,\qquad\text{with}\qquad
L=\int d^3\xx\,\Lag\bigl(\Phi(\xx,t),\dot\Phi(\xx,t),\nabla\Phi(\xx,t)\bigr)\,,
\]
i.e.\ $L$ is a $d^3\xx$ integral of a \emph{spatial density} $\Lag$, which
depends on $\Phi$, $\dot\Phi$, and $\nabla\Phi$, all taken \emph{at the same
point} $\xx$. This means that the field degrees of freedom $\Phi(\xx)$ at
some point $\xx$ do not have a direct interaction with $\Phi(\xx')$ located
at some other point $\xx'$, with $|\xx-\xx'|$ greater than zero. Allowing
such interactions would lead to possible terms in $L$ like
\[
L_{\text{nonlocal}}=\int F\bigl(\Phi(\xx,t),\Phi(\xx',t)\bigr)\,d^3\xx\,d^3\xx'
\]
(``nonlocal interactions''). Nonlocal interactions do not go well along with
the relativistic invariance. For if $S$ includes, say, double integrals over
$\xx$ (as above), relativistic invariance of $S$ forces one to include the
terms which are nonlocal in time as well, like
\[
\int F'\bigl(\Phi(t),\Phi(t')\bigr)\,dt\,dt'\,,
\]
which may violate \emph{causality}: The state in the future affects the
dynamics now. Locality and causality are deeply connected.

The above local form,
$L=\int\Lag\bigl(\Phi(\xx,t),\dot\Phi(\xx,t),\nabla\Phi(\xx,t)\bigr)\,d^3\xx$,
describes the dynamics in which the interaction occurs only between the
degrees of freedom associated with ``infinitely close'' spatial points
(through the $\nabla\Phi$ in $L$).

Also, I did not include higher spatial derivatives of $\Phi$ as the arguments
in $\Lag$. Although such dependence would not violate locality (on a
superficial level), including such dependencies in a relativistic theory
would require to add higher time derivatives, like $\ddot\Phi$, as well. Such
dependence is not allowed in conventional Lagrangian dynamics, which requires
that the time evolution $\Phi(\xx,t)$ is completely determined by the initial
values of $\Phi(\xx)$ and $\dot\Phi(\xx)$ at $t=t_0$.

Anyhow, we start with the above local form of the action. In classical field
theory the field equations of motion follow from the action principle
\[
\delta S=0\,.
\]
Let us combine the spatial and time variable into the 4-coordinate
$x=(x^\mu)=(x^0,\xx)$, with $x^0=t$, and write $\Phi(x)$ instead of
$\Phi(\xx,t)$; the Lagrangian density then is written as
\[
\Lag\bigl(\Phi(x),\partial_\mu\Phi(x)\bigr)\,,
\qquad\text{where}\qquad
\partial_\mu\equiv\frac{\partial}{\partial x^\mu}=\left(\frac{\partial}{\partial t},\nabla\right)\,.
\]
A simple exercise in functional analysis (look up the textbook, page 15 in
PS) yields the classical field equations of motion
\begin{equation}\label{eq:EL}
\partial_\mu\left(\frac{\partial\Lag}{\partial\bigl(\partial_\mu\Phi(x)\bigr)}\right)-\frac{\partial\Lag}{\partial\Phi(x)}=0\,.
\end{equation}

\begin{remark}
Essential step in deriving \eqref{eq:EL} is integrating by parts. The total
derivative term
\[
\int d^4x\,\frac{\partial}{\partial x^\mu}\left(\delta\Phi(x)\,\frac{\partial\Lag}{\partial(\partial_\mu\Phi(x))}\right)
\]
can be reduced to the integral over the boundary (Gauss theorem), which
depends only on the variations $\delta\Phi(x)$ at the boundary; this term does
not affect the ``bulk'' equations of motion \eqref{eq:EL}, which hold at each
point $x$ of the space-time. But if the field dynamics is considered in a
space-time domain $D$, with the boundary $\partial D$, the above term
contributes to the variation of $S$ via the variations of the boundary values
$\delta\Phi(x)\,|_{x\in\partial D}$ of the fields,
\[
\int_D d^4x\,\partial_\mu(\dots)=\int_{\partial D}\delta\Phi(x)\,\frac{\partial\Lag}{\partial\,\partial_\mu\Phi(x)}\,d\Sigma_\mu(x)\,,
\]
where $d\Sigma_\mu(x)$ is the normal\footnote{``Normal'' means that
$d\Sigma_\mu(x)\,T^\mu(x)=0$ for any $T^\mu(x)$ tangent to the boundary
$\partial D$, at any point $x\in\partial D$.} element of the boundary
hypersurface $\partial D$. In this situation the total derivative term
contributes to the boundary conditions (see the homework
Problem~\ref{hw:1}). The default setup is $D=$ full infinite space-time, with
the condition that $\Phi(x)$ decays sufficiently fast at both space and time
infinities (``scattering asymptotic conditions''); in this setup the above
derivative term can be ignored.
\end{remark}

The task of the classical field theory is to find solutions of these
equations with such and such initial or boundary conditions.

One of the most important general results about classical field theory is
the relation between symmetries and conservation laws known as the

\subsection*{Noether's theorem (pp.~17--18 of PS)}

What is a symmetry? Suppose we have a family of continuous transformations of
the field variables
\[
\Phi(x)\to\Phi'(x)=F\bigl(x,\Phi(x)\bigr)
\]
such that the action $S$ does not change
\[
S[\Phi]=S[\Phi']\,.
\]
The term continuous transformations means here that they depend on some
continuous parameter(s), $F=F_s$, in such a way that
\[
F_0\bigl(x,\Phi(x)\bigr)=\Phi(x)\qquad\text{-- no transformation.}
\]
(discrete symmetries is a separate story). Then we can take an infinitesimal
transformation
\[
\Phi(x)\to\Phi'(x)=\Phi(x)+\epsilon\,E\bigl(x,\Phi(x)\bigr)\,,
\qquad\text{where}\qquad
E=\frac{d}{ds}F\Big|_{s=0}\,,
\]
with an infinitesimal $\epsilon$. Substituting in the Lagrangian density, one
finds
\[
\Lag(\Phi',\partial_\mu\Phi')=\Lag(\Phi,\partial_\mu\Phi)
+\epsilon\left[\frac{\partial\Lag}{\partial\Phi}\,E(x,\Phi)+\frac{\partial\Lag}{\partial(\partial_\mu\Phi)}\,\partial_\mu E(x,\Phi)\right]+O(\epsilon^2)\,.
\]

The invariance (i.e.\ the requirement that the action $S=\int\Lag\,d^4x$
does not change) implies
\[
\bigl[\cdots\bigr]=\partial_\mu\tilde J^\mu(\Phi,\partial_\mu\Phi,x)\,,
\]
with some $\tilde J^\mu$ which depend on our field variables in a local way
(it may also have separate dependence on $x$ if $E$ does). Transforming the
left-hand side by parts one finds
\[
\partial_\mu\left(\frac{\partial\Lag}{\partial(\partial_\mu\Phi)}\,E(x,\Phi)\right)
+\left[\frac{\partial\Lag}{\partial\Phi}-\partial_\mu\left(\frac{\partial\Lag}{\partial(\partial_\mu\Phi)}\right)\right]E(x,\Phi)=\partial_\mu\tilde J^\mu\,.
\]
Now, suppose that $\Phi=\Phi_c$, a solution of the classical field equations
of motion. Then the second term vanishes, and we obtain the \emph{continuity
equation}
\[
\partial_\mu J^\mu=0\,,
\qquad\text{where}\qquad
J^\mu(\Phi,\partial\Phi,x)=\frac{\partial\Lag}{\partial(\partial_\mu\Phi)}\,E(x,\Phi)-\tilde J^\mu(\Phi,\partial\Phi,x)\,,
\]
which holds as long as $\Phi$ solves the classical field equations.

The continuity equation then implies a conservation law. Indeed, by the Gauss
theorem
\[
\oint_\Sigma J^\mu\,d\Sigma_\mu=\int_D\partial_\mu J^\mu\,d^4x=0\,.
\]
Here $\Sigma=\partial D$ is any closed 3-dimensional hypersurface in the
4-dimensional space-time, and
$d\Sigma_\mu=\frac{1}{3!}\,\epsilon_{\mu\nu\lambda\rho}\,dx^\nu dx^\lambda dx^\rho$
is the normal element of this hypersurface. Taking the hypersurface to be a
big slab between two equal-time hyperplanes and bringing the spatial boundary
to the spatial infinity (figure~\ref{fig:slab}),
\begin{figure}[ht]
\centering\figfile{slab}
\caption{The closed hypersurface $\Sigma=\partial D$: a slab bounded by the
equal-time hyperplanes $t=t_1$ and $t=t_2$, with the side boundary taken to
spatial infinity. The charges $Q_{t_1}$ and $Q_{t_2}$ are the fluxes of $J^\mu$
through the bottom and top faces.}
\label{fig:slab}
\end{figure}
and also assuming that all fields decay at $\xx\to\infty$, one finds
\[
Q_{t_1}=Q_{t_2}\,,\qquad\text{where}\qquad Q_t=\int J^0\,d^3\xx\,,
\]
that is $Q$ is an integral of motion for the classical field equations. Let
me stress again that both the continuity equation $\partial_\mu J^\mu=0$ and
the conservation law $\frac{d}{dt}Q=0$ are satisfied only if we take
$\Phi=\Phi_c$ (in terminology of QFT they are said to hold ``on shell'').

\subsection*{Energy-Momentum}

The most common symmetry we encounter in field theory is the translational
symmetry. Assume that the Lagrangian density $\Lag(\Phi,\partial\Phi)$ has no
$x$ dependence except for that coming through the field variable $\Phi(x)$.
It means that the dynamics looks the same if we shift the space-time
coordinates
\[
x\to x'=x+a
\]
by a constant 4-vector $a$. This can be viewed as the symmetry with respect
to the following transformation of the field variables
\[
\Phi(x)\to\Phi'(x)=\Phi(x+a)\,,
\]
which leaves the action invariant. This field transformation has the
infinitesimal form
\[
\Phi'(x)=\Phi(x)+da^\mu\,\partial_\mu\Phi(x)\,,
\]
i.e.\ in this example $E(\Phi,\partial\Phi)=\partial_\mu\Phi$. When we do
this transformation the action is going to stay invariant, but the Lagrangian
density changes by the total derivative
\[
\Lag'=\Lag+da^\nu\,\partial_\nu\Lag\,,
\]
where the derivative in the last term can be written as
$\partial_\mu(\delta^\mu_\nu\,\Lag)$, and has the same meaning as
$\partial_\mu\tilde J^\mu$ in the general argument above.

Note that there are four independent symmetries associated with four
components of $a^\mu$. Correspondingly, there are four Noether's currents
\[
T^\mu{}_\nu=-\frac{\partial\Lag}{\partial(\partial_\mu\Phi)}\,\partial_\nu\Phi+\delta^\mu_\nu\,\Lag\,,
\]
which therefore satisfy the continuity equations
\[
\partial_\mu T^\mu{}_\nu=0\qquad\text{on shell.}
\]
This object is known as the \textbf{energy-momentum tensor}. It leads to four
conserved charges
\[
E=-\int T^0{}_0\,d^3\xx\,,
\qquad\text{and}\qquad
P^i=\int T^0{}_i\,d^3\xx\,,\quad\text{with}\quad i=1,2,3\,.
\]
Of these, $E$ is interpreted as the energy of the field configuration
$\Phi(\xx,t)$, and $\mathbf{P}$ is interpreted as its momentum. Therefore one
may take
\begin{align*}
T^{00}&=-T^0{}_0=\mathcal{E}&&\text{for the energy density,}\\
T^{0i}&=T^0{}_i=\mathcal{P}^i&&\text{for the momentum density.}
\end{align*}
The spatial components $T^{ij}$ constitute the (spatial) \emph{stress tensor}.
It is interpreted as usual: Consider any domain $D=D^3$ in the 3-dimensional
space (at a given time). The momentum of the field inside this domain is
\[
P^i_D=\int_D\mathcal{P}^i(\xx)\,d^3\xx\,.
\]
According to the momentum theorem (Newton's second law),
\[
\frac{d}{dt}P^i_D=F^i_D\,,\qquad\text{the net force applied to the field inside }D\,.
\]
Since the interaction is local, this force must come from the surface forces,
\[
F^i_D=-\int_{\partial D}T^{ij}(\xx)\,d\sigma_j=-\int_D\partial_jT^{ij}(\xx)\,d^3\xx\,.
\]
Thus, $dF^i=T^{ij}(\xx)\,d\sigma_j$ is the force exerted by the field on the
surface element $d\sigma_j$.

\subsection*{Klein-Gordon Field}

Let us consider a simple field theory known as the \emph{Klein-Gordon} theory
(it will be our main example for some time). It involves a single-component
scalar field $\Phi(x)=\varphi(\xx,t)$, and the action is
\[
S=\int d^4x\,\left[-\frac12\,\partial_\mu\varphi\,\partial^\mu\varphi-\frac{m^2}{2}\,\varphi^2\right]
\equiv\int dt\,d^3\xx\,\left[\frac12\,\dot\varphi^2-\frac12\,(\nabla\varphi)^2-\frac{m^2}{2}\,\varphi^2\right].
\]
This form of the action shows that the theory is Lorentz-invariant, and the
field $\varphi$ is a scalar (i.e.\ $\varphi(x)$ has definite value at each
space-time point $x$, independent of the choice of the inertial frame).
Variation yields the Klein-Gordon field equation
\[
-\partial_\mu\partial^\mu\varphi+m^2\,\varphi=\partial_t^2\varphi-\nabla^2\varphi+m^2\,\varphi=0\,.
\]
Applying the above equations one finds for the energy-momentum tensor
\[
T_{\mu\nu}=\partial_\mu\varphi\,\partial_\nu\varphi-\frac12\,\eta_{\mu\nu}\,\bigl(\partial_\lambda\varphi\,\partial^\lambda\varphi+m^2\,\varphi^2\bigr)\,,
\]
with $\eta_{\mu\nu}=\operatorname{diag}(-1,1,1,1)$, so that
$\mathcal{E}=T^{00}=-T^0{}_0$ and $\mathcal{P}^i=T^{0i}=T^0{}_i$. Thus,
\[
\mathcal{E}=\frac12\Bigl(\dot\varphi^2+(\nabla\varphi)^2+m^2\,\varphi^2\Bigr)\,,
\qquad\text{and}\qquad
\boldsymbol{\mathcal{P}}=-\dot\varphi\,\nabla\varphi\,.
\]
We see in this example that the energy-momentum tensor $T_{\mu\nu}$ is
symmetric,
\[
T_{\mu\nu}=T_{\nu\mu}\,.
\]

\begin{exercise}
Show that in any Lorentz-invariant theory involving only a scalar field
$\varphi(x)$ the energy-momentum tensor always comes out symmetric,
$T_{\mu\nu}=T_{\nu\mu}$. Check that in such theory all components of the
antisymmetric tensor
\[
M^{\mu\nu}=\int\bigl[x^\mu\,T^{\nu0}-x^\nu\,T^{\mu0}\bigr]\,d^3\xx
\]
are conserved on-shell.
\end{exercise}

The equations of motion of the Klein-Gordon theory
\[
-\partial_\mu\partial^\mu\varphi+m^2\,\varphi=\partial_t^2\varphi-\nabla^2\varphi+m^2\,\varphi=0
\]
are \emph{linear}. That is, any linear combination of solutions is again a
solution. (In general, this property defines the \emph{free field theory}).
Therefore, general solution is a superposition of ``elementary solutions''.
In the infinite space-time, one can take the plane waves
\[
e^{ipx}=e^{-ip^0x^0+i\pp\cdot\xx}\qquad\text{with}\qquad p^0=\omega_\pp:=\sqrt{\pp^2+m^2}
\]
as well as the complex-conjugate waves
\begin{equation}\label{eq:KGgeneral}
\varphi(x)=\int\frac{d^3\pp}{2\omega_\pp\,(2\pi)^3}\,\bigl(A_\pp\,e^{ipx}+A^*_\pp\,e^{-ipx}\bigr)\,.
\end{equation}
This form represents general solution of the Klein-Gordon equation, which is
bounded at the spatial infinity.

\begin{hwproblem}{1}\label{hw:1}
Consider a classical field theory in the semi-infinite space,
$x^1,x^2\in(-\infty,\infty)$ but $x^3\in[0,\infty)$, with the boundary at
$x^3=0$. Assume that there is a single scalar field $\varphi(\xx,t)$, and the
Lagrangian has the form
\[
L=\int_0^\infty dz\int d^2x_\parallel\,\left[\frac12\,\dot\varphi^2-\frac12\,\bigl(\nabla\varphi\bigr)^2-\frac{m^2}{2}\,\varphi^2\right]
-\int d^2x_\parallel\,V\bigl(\varphi_B(x_\parallel,t)\bigr)\,,
\]
with some function $V(\varphi_B)$. Here $z\equiv x^3$, and
$x_\parallel=(x^1,x^2)$ are the coordinates along the boundary $z=0$, and
$\varphi_B(x_\parallel,t)$ stands for the boundary values of $\varphi$,
\[
\varphi_B(x_\parallel,t)=\varphi(\xx,t)\big|_{z=0}\,.
\]
\begin{hwparts}
\item Find the field equations of motion, including the boundary conditions
at $z=0$.
\item Which of the quantities $E$, $P^1$, $P^2$, $P^3$ are still conserved?
\item Assuming that $V(\varphi_B)=M\varphi_B^2$ with $M>0$, find general
solutions of the field equations, bounded at $z\to\infty$.
\item Consider now $V(\varphi_B)=M\varphi_B^2$ with $M<0$. Find additional
solutions which appear in this case.
\end{hwparts}
\end{hwproblem}
